% ==============================================================================
% SLIDES - SEMANA 12: FRAMEWORKS DE GOVERNANÇA (ISO 38500, COBIT, ITIL)
% ==============================================================================

% ------------------------------------------------------------------------------
% 1. SLIDE DE CAPA
% ------------------------------------------------------------------------------
\senaicover
  {Frameworks de Governança \& Serviços}
  {ISO/IEC 38500, COBIT 2019, ITIL v4 e Seleção Justificada de Processos}
  {Unidade Curricular: Governança de TI}
  {Prof. André Souza}

% ------------------------------------------------------------------------------
% 2. VISÃO GERAL / AGENDA
% ------------------------------------------------------------------------------
\begin{senaiframe}{Visão Geral \& Agenda}{Os 4 Eixos da Aula}

  \begin{columns}[t]
    \begin{column}{0.48\textwidth}
      \begin{tcolorbox}[senaibluecard, title={1. ISO/IEC 38500}]
        \begin{itemize}
          \item Os 6 Princípios da Governança de TI.
          \item O ciclo EDM (Avaliar, Dirigir, Monitorar).
        \end{itemize}
      \end{tcolorbox}
      
      \vspace{4pt}
      
      \begin{tcolorbox}[senaiambercard, title={3. ITIL v4}]
        \begin{itemize}
          \item Sistema de Valor de Serviço (SVS).
          \item As 34 Práticas de Gestão de Serviços.
        \end{itemize}
      \end{tcolorbox}
    \end{column}

    \begin{column}{0.48\textwidth}
      \begin{tcolorbox}[senairedcard, title={2. COBIT 2019}]
        \begin{itemize}
          \item Estrutura dos 40 Objetivos (EDM, APO, BAI, DSS, MEA).
          \item Fatores de Design e Governança Sob Medida.
        \end{itemize}
      \end{tcolorbox}
      
      \vspace{4pt}
      
      \begin{tcolorbox}[senaigreencard, title={4. Capítulo 7 do PDGTI}]
        \begin{itemize}
          \item Seleção e Justificativa dos Frameworks.
          \item Matriz de Intersecção dos Modelos.
        \end{itemize}
      \end{tcolorbox}
    \end{column}
  \end{columns}

\end{senaiframe}

% ------------------------------------------------------------------------------
% 3. ISO 38500 & COBIT 2019
% ------------------------------------------------------------------------------
\begin{senaiframe}{Módulo 1 • Diretrizes \& Governança}{ISO/IEC 38500 \& COBIT 2019}

  \begin{columns}[t]
    \begin{column}{0.48\textwidth}
      \begin{tcolorbox}[senaibluecard, title={ISO/IEC 38500 (Nível Conselho)}]
        \begin{itemize}
          \item \textbf{6 Princípios:} Responsabilidade, Estratégia, Aquisição, Desempenho, Conformidade, Comportamento.
          \item \textbf{Modelo EDM:} Evaluate, Direct, Monitor.
        \end{itemize}
      \end{tcolorbox}
    \end{column}

    \begin{column}{0.48\textwidth}
      \begin{tcolorbox}[senairedcard, title={COBIT 2019 (Visão Holística)}]
        \begin{itemize}
          \item \textbf{5 Domínios:} EDM (Governança) + APO, BAI, DSS, MEA (Gestão).
          \item \textbf{40 Objetivos:} Seleção dos processos prioritários baseados nos Fatores de Design.
        \end{itemize}
      \end{tcolorbox}
    \end{column}
  \end{columns}

\end{senaiframe}

% ------------------------------------------------------------------------------
% 4. ITIL V4 & COMPLEMENTARIDADE
% ------------------------------------------------------------------------------
\begin{senaiframe}{Módulo 2 • Serviços \& Integração}{ITIL v4 \& Hierarquia de Abstração}

  \begin{columns}[t]
    \begin{column}{0.48\textwidth}
      \begin{tcolorbox}[senaiambercard, title={ITIL v4 (Gestão de Serviços)}]
        \begin{itemize}
          \item \textbf{SVS:} Sistema de Valor de Serviço e Cadeia de Valor.
          \item \textbf{7 Princípios:} Foco no valor, começar de onde está, manter simples e prático.
        \end{itemize}
      \end{tcolorbox}
    \end{column}

    \begin{column}{0.48\textwidth}
      \begin{tcolorbox}[senaicard, title={Hierarquia dos Frameworks}]
        \begin{itemize}
          \item \textbf{ISO 38500:} *O Quê* fazer no Conselho.
          \item \textbf{COBIT 2019:} *Como* governar e estruturar.
          \item \textbf{ITIL v4:} *Como* executar os serviços diariamente.
        \end{itemize}
      \end{tcolorbox}
    \end{column}
  \end{columns}

\end{senaiframe}

% ------------------------------------------------------------------------------
% 5. REFERÊNCIAS BIBLIOGRÁFICAS
% ------------------------------------------------------------------------------
\begin{senaiframe}{Referências Bibliográficas}{Fontes \& Leitura Recomendada}

  \begin{tcolorbox}[senaicard, title={Obras de Referência}]
    \begin{itemize}
      \item \textbf{ISO/IEC.} *ISO/IEC 38500: Governance of IT for the organization.* ISO, 2015.
      \item \textbf{ISACA.} *COBIT 2019 Framework: Governance and Management Objectives.* Schaumburg: ISACA, 2018.
      \item \textbf{AXELOS.} *ITIL Foundation: ITIL 4 Edition.* London: TSO, 2019.
    \end{itemize}
  \end{tcolorbox}

\end{senaiframe}

% ------------------------------------------------------------------------------
% 6. APLICAÇÃO NO PDGTI
% ------------------------------------------------------------------------------
\begin{senaiframe}{Aplicação no PDGTI}{Desenvolvimento do Capítulo 7}

  \vspace{0.2cm}
  \begin{tcolorbox}[senaigreencard, title={Capítulo 7 – Seleção e Justificativa de Frameworks}]
    \begin{itemize}
      \item \textbf{7.1 ISO 38500:} Aplicação do ciclo EDM no Comitê de TI da empresa.
      \item \textbf{7.2 COBIT 2019:} Seleção de no mínimo 3 Objetivos prioritários.
      \item \textbf{7.3 ITIL v4:} Seleção e justificativa das práticas de serviços (Incidentes, Mudanças, SLA).
      \item \textbf{7.4 Matriz de Intersecção:} Demonstrar a complementaridade prática entre os modelos.
    \end{itemize}
  \end{tcolorbox}

\end{senaiframe}
